\documentclass[11pt,a4paper]{article}

\usepackage[a4paper,top=25mm,bottom=27mm,left=27mm,right=27mm]{geometry}
\usepackage{fontspec}
\setmainfont{Times New Roman}
\setsansfont{Arial}
\setmonofont{Menlo}
\usepackage{microtype}
\usepackage{amsmath,amssymb,mathtools}
\usepackage{booktabs,longtable,array,tabularx}
\usepackage{float}
\usepackage{enumitem}
\usepackage{xcolor}
\usepackage{titlesec}
\usepackage{fancyhdr}
\usepackage{hyperref}
\usepackage{bookmark}
\usepackage{setspace}

\definecolor{SIELblue}{HTML}{173A5E}
\definecolor{SIELgray}{HTML}{5A6570}
\definecolor{SIELlight}{HTML}{EDF2F6}
\hypersetup{
  colorlinks=true,
  linkcolor=SIELblue,
  citecolor=SIELblue,
  urlcolor=SIELblue,
  pdftitle={Subjectivity-Intersection Mathematics: Perspective Intersection, Emergent O3, and Self-Re-entry},
  pdfauthor={Satoru Watanabe},
  pdfsubject={Public working preprint},
  pdfkeywords={subjectivity-intersection mathematics, perspective intersection, third perspective, O3, nonseparability, reciprocal learning, self-re-entry}
}

\titleformat{\section}{\Large\sffamily\bfseries\color{SIELblue}}{\thesection}{0.7em}{}
\titleformat{\subsection}{\large\sffamily\bfseries\color{SIELblue}}{\thesubsection}{0.7em}{}
\titleformat{\subsubsection}{\normalsize\sffamily\bfseries\color{SIELgray}}{\thesubsubsection}{0.7em}{}
\titlespacing*{\section}{0pt}{2.2ex plus .5ex}{1.1ex}
\titlespacing*{\subsection}{0pt}{1.8ex plus .4ex}{0.8ex}

\pagestyle{fancy}
\fancyhf{}
\fancyhead[L]{\small\sffamily Subjectivity-Intersection Mathematics}
\fancyhead[R]{\small\sffamily Public Working Preprint v3.0 Draft}
\fancyfoot[C]{\small\thepage}
\renewcommand{\headrulewidth}{0.3pt}
\setlength{\headheight}{14pt}
\setlength{\parindent}{1.2em}
\setlength{\parskip}{0.35em}
\setlength{\emergencystretch}{2em}
\setlist{itemsep=0.25em,topsep=0.45em}
\renewcommand{\arraystretch}{1.16}

\newcommand{\OA}{\ensuremath{O_{A}}}
\newcommand{\OB}{\ensuremath{O_{B}}}
\newcommand{\Othree}{\ensuremath{O_{3}}}
\newcommand{\statussupported}{\textsf{SUPPORTED}}
\newcommand{\statusnotsupported}{\textsf{NOT SUPPORTED}}

\begin{document}

\begin{titlepage}
\thispagestyle{empty}
\vspace*{18mm}
{\sffamily\bfseries\color{SIELblue}\fontsize{24}{29}\selectfont
\mbox{Subjectivity-Intersection}\\Mathematics\par}
\vspace{6mm}
{\sffamily\fontsize{18}{23}\selectfont
Perspective Intersection, Emergent O3,\\and Self-Re-entry\par}
\vspace{7mm}
{\sffamily\fontsize{15}{19}\selectfont
A Staged Preregistered Computational\\Research Program\par}

\vspace{18mm}
{\large Satoru Watanabe\par}
\vspace{2mm}
{\normalsize Subjectivity Intersection Emergence Lab\par}
{\normalsize Correspondence: \href{mailto:satoru@siel.global}{satoru@siel.global}\par}

\vfill
\noindent\colorbox{SIELlight}{%
  \parbox{0.93\textwidth}{%
  \vspace{2mm}
  \textbf{Document status:} Public Working Preprint, Version 3.0 Draft\\
  \textbf{Date:} 5 August 2026\\
  \textbf{DOI (reserved):} \href{https://doi.org/10.5281/zenodo.21796460}{10.5281/zenodo.21796460}\\
  \textbf{Previous citable version:} \href{https://doi.org/10.5281/zenodo.21791244}{doi:10.5281/zenodo.21791244}\\
  \textbf{Repository:} \href{https://github.com/SIEL-Research/Subjectivity-Intersection-Mathematics}{SIEL-Research/Subjectivity-Intersection-Mathematics}\\
  Version 3 reorganizes the program around reciprocal perspective intersection. It identifies the learned nonseparable component as a distributed operational third perspective, O3, and treats relational generation as a consequence of that intersection.
  \vspace{2mm}
  }}

\vspace{10mm}
{\small Copyright \textcopyright\ 2026 Satoru Watanabe. All rights reserved.\par}
\end{titlepage}

\begin{abstract}
What happens when two distinct perspectives do not merely exchange information but each takes the other into its own state formation? This paper develops \emph{Subjectivity-Intersection Mathematics} as a staged preregistered computational program centered on reciprocal perspective intersection. A perspective is defined operationally as a history-dependent state organization that transforms received information and constrains subsequent state and output formation. Two perspectives, \OA\ and \OB, intersect when each is recurrently changed through the other while their distinction is preserved.

The central hypothesis is that reciprocal intersection can generate a third perspective, \Othree, that is reducible to neither participant's perspective, is not installed as a third agent or state register, and subsequently constrains both participants and their next intersection. The computational realization begins with two independent persistent states trained only on delayed reciprocal recall. It contains no state named C or O3, no pair identifier, carrier target, teacher label, or carrier-specific loss. After learning, counterfactual comparison extracts a distributed joint component C. Selective erasure, cross-pair exchange, signal-free transport, exact reconstruction, and equal-capacity disconnected controls test its identity and causal role.

Across the staged sequence, Experiment 005 found an emergent interaction-specific solution class in 84 of 96 registered architecture-seed units; the independent control produced no corresponding component. Experiment 006 preserved a preregistered failure despite transport in all 96 units. Experiment 006R then separated functional self-re-entry from external identifiability and supported all 13 revised primary conditions on new seeds, with all-three-transition transport in 190 of 192 units. Experiment 006A directly tested whether the distributed result was only two additive directed memories. All ten preregistered conditions passed across 48 new seeds and five exactly capacity-matched architectures. Four interacting architectures produced nonzero, bilateral, transported, and exchange-sensitive joint terms, whereas a fully task-competent disconnected dual relay had maximum joint norm $2.93\times10^{-17}$.

Within the stated operational definition, these results support identifying C as an emergent distributed third perspective, \Othree. Its two directed realizations are located within A and B, but their function is jointly nonseparable: neither side alone nor two disconnected reciprocal memories reproduce it. Removing present \Othree\ alters both next participant states and the next extracted intersection. The experimentally supported computational grammar is therefore
\[
  (\OA,\OB) \xrightarrow{\text{reciprocal intersection}} \Othree
  \xrightarrow{\text{self-re-entry}} (\OA',\OB',\Othree').
\]
The result establishes an operational third perspective in the tested model class; claims about phenomenal experience or independent subjecthood require additional criteria beyond the present computation.
\end{abstract}

\noindent\textbf{Keywords:} subjectivity-intersection mathematics; perspective intersection; third perspective; O3; nonseparability; reciprocal learning; distributed realization; self-re-entry; preregistration

\section*{Principal contributions}
\addcontentsline{toc}{section}{Principal contributions}
{\small
\begin{enumerate}[itemsep=0.1em,topsep=0.25em]
  \item The primary mathematical object is reformulated as \emph{perspective intersection}, not relation in the abstract.
  \item An operational perspective is defined by history-dependent transformation and constraint on subsequent state and output formation.
  \item C is identified as a distributed operational third perspective, \Othree, when it emerges from reciprocal intersection and satisfies non-reduction, joint nonseparability, pair specificity, bilateral action, transport, and self-re-entry.
  \item Experiment 006A distinguishes one nonseparable O3 contribution from two equally competent but disconnected directed memories.
  \item The staged evidence supports the closure $(\OA,\OB)\rightarrow\Othree\rightarrow(\OA',\OB',\Othree')$ without an initial O3 state, O3 target, or O3-specific objective.
  \item The original negative decision of Experiment 006 is retained and used to distinguish functional self-re-entry from external identifiability.
\end{enumerate}
}

\clearpage
{\small
\setlength{\parskip}{0pt}
\tableofcontents
}
\clearpage

\section{Introduction: from relation to intersection}

Interactions are commonly described through messages, shared variables, latent edges, coordination patterns, or models one agent holds of another. These descriptions begin with participants and characterize what passes between them. Subjectivity-Intersection Mathematics asks a different question: can the reciprocal taking-up of one perspective by another generate a perspective that belongs to neither perspective separately and then participate in the formation of both?

The distinction matters. A relation may exist whenever A affects B. A perspective intersection requires two differentiated perspectives and reciprocal transformation: A receives B through A's own state organization, while B receives A through B's own. The two directed transformations may remain separate. The stronger possibility is that they become jointly nonseparable and function as a third perspective.

Let \OA\ and \OB\ denote the operational perspectives of A and B. The canonical structure of the program is
\begin{equation}
 (\OA,\OB)
 \xrightarrow{\mathcal X}
 \Othree
 \xrightarrow{\mathcal E}
 (\OA',\OB',\Othree'),
 \label{eq:canonical}
\end{equation}
where $\mathcal X$ is reciprocal intersection and $\mathcal E$ is self-re-entry through unchanged dynamics. The first arrow is not averaging, fusion, consensus, or the installation of a shared memory. The second is not an external feedback channel added after the fact. It states that the generated perspective becomes part of the conditions under which the next two perspectives and their next intersection are formed.

This paper reports a sequence of preregistered computational tests of that grammar. Early experiments constructed an explicit carrier to establish operational sufficiency. Later experiments removed the carrier from the architecture and trained only an ordinary reciprocal task. The final stages asked whether the emergent distributed result was causal, temporally recurrent, architecture-general, pair-specific, and one jointly nonseparable contribution rather than two directed memories.

\section{A mathematical definition of perspective intersection}

\subsection{Operational perspective}

A perspective is not identified with an input token or stored fact. It is the state-dependent way in which a system receives, transforms, and continues from what is encountered. For participant $i\in\{A,B\}$, define
\begin{equation}
 O_i(t)=\bigl(h_i(t),\mathcal H_i(t),F_i,D_i\bigr),
 \label{eq:perspective}
\end{equation}
where $h_i$ is the current persistent state, $\mathcal H_i$ its history, $F_i$ its update organization, and $D_i$ its output map. Two systems can receive the same symbol yet instantiate different perspectives because their histories and transformations differ.

In the reciprocal-recall experiments, discrete symbols are minimal probes, not perspectives themselves. The operational perspective is the recurrent organization through which A incorporates B-originating information and through which B incorporates A-originating information.

\subsection{Reciprocal intersection}

Let the coupled updates be
\begin{align}
 h_A(t+1)&=F_A(h_A(t),h_B(t),x_A(t)),\\
 h_B(t+1)&=F_B(h_B(t),h_A(t),x_B(t)).
 \label{eq:coupled}
\end{align}
Reciprocal intersection occurs when both cross-dependencies are active and alter subsequent state formation. It preserves distinction: $h_A$ and $h_B$ remain separate persistent states. Intersection is therefore not merger but mutually conditioned difference.

The two directed traces may be written
\begin{align}
 R_t^{A\leftarrow B}&=h_A^{AB}(t)-h_A^{A\text{-only}}(t),\\
 R_t^{B\leftarrow A}&=h_B^{AB}(t)-h_B^{B\text{-only}}(t).
 \label{eq:directed}
\end{align}
These traces show that each perspective has taken up the other. They do not yet establish a third perspective. That requires evidence that the two directed realizations participate in one nonseparable operation.

\subsection{The third perspective O3}

Define the joint post-learning intersection term
\begin{equation}
 C_t=H_t(ab)-H_t(a0)-H_t(0b)+H_t(00),
 \label{eq:joint}
\end{equation}
where all four passes use identical weights, matched inputs, times, initial conditions, and evaluation episodes. The term removes the two one-sided histories and restores their common baseline. It is zero for a purely additive pair of disconnected directed memories.

Within this paper, C is identified as an operational third perspective \Othree\ when the following conditions hold:
\begin{enumerate}
  \item it is absent as a named state, target, label, or specialized loss before learning;
  \item it appears only under reciprocal interaction and post-learning counterfactual extraction;
  \item it is reducible to neither \OA\ nor \OB\ under the frozen decomposition;
  \item its A-side and B-side realizations are jointly nonseparable against an equally competent additive control;
  \item its selective deletion changes both receivers;
  \item its replacement by an equal-norm component from another pair loses the original function;
  \item it contributes through unchanged dynamics to the next participant states and next extracted intersection.
\end{enumerate}

This definition does not require a central O3 register. O3 can be distributed across A and B while remaining one operation. Its unity is functional and counterfactual: the two locations cannot be separated without losing the measured joint effect.

\subsection{Self-re-entry}

Let $X_t=(h_A(t),h_B(t))$ and let $\mathcal I_{-C_t}$ remove only the extracted O3 contribution. Under unchanged update $F$,
\begin{align}
 X_{t+1}&=F(X_t),\\
 X_{t+1}^{(-C)}&=F(\mathcal I_{-C_t}(X_t)),\\
 \Delta X_{t+1}^{C}&=X_{t+1}-X_{t+1}^{(-C)}.
 \label{eq:state-return}
\end{align}
Applying the frozen extractor $\mathcal R$ at the next state gives
\begin{equation}
 \Delta C_{t+1}^{C}
 =\mathcal R(F(X_t))-\mathcal R(F(\mathcal I_{-C_t}(X_t))).
 \label{eq:o3-return}
\end{equation}
Equation~\eqref{eq:state-return} tests return into A and B; Equation~\eqref{eq:o3-return} tests return into the next intersection. Exact reconstruction compares the observed missing contribution with the contribution predicted through the unchanged update. Thus O3 is not merely retained information. It is a current constraint on what the next perspectives and next intersection can become.

\section{Relation to neighboring concepts}

\subsection{A's model of B and B's model of A}

A representation held by A about B belongs locally to A; the converse belongs locally to B. The present claim begins only after both directed traces have been isolated. Experiment 006A asks whether those traces remain additive. The disconnected dual relay demonstrates the additive case: both agents solve reciprocal recall, yet Equation~\eqref{eq:joint} cancels to numerical zero. O3 denotes the additional jointly nonseparable organization found in interacting systems.

\subsection{Shared memory and third agent}

A shared memory localizes common information in an explicitly accessible place. A third agent possesses its own independently installed state and update. Neither is present in Experiments 005--006A. The emergent O3 is reconstructed from distributed changes inside A and B after learning.

\subsection{Theory of mind, communication, and relational models}

Theory-of-mind models normally infer beliefs or goals attributed to another agent. Emergent-communication models learn messages or protocols. Graph and relational models learn edges or relation types. These are adjacent but distinct objects. O3 is defined by reciprocal formation, nonseparability, pair specificity, bilateral causal action, and self-re-entry, not by semantic attribution, message identity, or edge prediction alone.

\subsection{Functional perspective and phenomenal subjectivity}

The term \emph{perspective} is used here in a mathematical-functional sense fixed by Equation~\eqref{eq:perspective}. The experiment establishes that the extracted C satisfies the registered structural and causal conditions used to identify an operational O3 in this model class. Phenomenal experience and independent subjecthood require additional bridge criteria and are not variables in the present experiment.

\subsection{Connection to Subjectivity Intersection Ontology}

The companion ontology preprint, \emph{Subjectivity Intersection Ontology: A Comparative Ontological Framework for Third Subjectivity}~\cite{watanabe2026sio}, defines subjectivity as perspective and distinguishes two related structures. Its ontological prototype is a simultaneous triadic closure of Absolute Subjectivity, Relative Subjectivity, and Intersection Subjectivity. For the intersection of multiple Relative Subjectivities, it separately gives the operational form $AB\rightarrow C\rightarrow AB$: differentiated perspectives generate an irreducible third perspective, and that third acts back upon the perspectives through which it arose.

The present program connects directly to the second structure. Its \OA\ and \OB\ are finite operational perspectives rather than claims about Absolute or Relative Subjectivity as such. Their reciprocal intersection generates a jointly nonseparable C; Experiment 006A distinguishes this C from two additive directed memories, and Experiments 006 and 006R test whether present C constrains later A, B, and C formation. The computational sequence
\[
 (\OA,\OB)\xrightarrow{\text{intersection}}\Othree
 \xrightarrow{\text{self-re-entry}}(\OA',\OB',\Othree')
\]
therefore supplies an executable realization of the ontology's multiple-relative-subject operational grammar. The points of connection are differentiated perspectives, generation through intersection, irreducibility of the third term, bilateral return, and closure through renewed intersection.

The ontological prototype and the computational realization nevertheless remain different in kind. The ontology begins from primordial unity, treats subjectivity as ontological perspective, makes the constitution of its triad simultaneous, and proposes recurrence across levels. The experiment begins from two separately initialized finite systems, represents their histories numerically, observes a temporal sequence of state updates, and tests a counterfactually extracted contribution. It does not test primordial unity, the projection of Absolute into Relative Subjectivity, phenomenal subjecthood, or cross-level recurrence. Nor does temporal self-re-entry by itself establish the ontology's simultaneous constitutive closure.

The connection is therefore neither an identity claim nor a merely decorative analogy. It is a directional bridge. The ontology specifies the structural hypothesis and distinguishes its ontological prototype from its operational expression; the mathematics turns the operational expression into conditions that can succeed or fail. Conversely, the experiments constrain which computational organizations can legitimately model $AB\rightarrow C\rightarrow AB$. Within that bridge, the present C is an \emph{operational third perspective}; its identification with Intersection Subjectivity remains a further bridge hypothesis requiring criteria that connect functional perspective to subjectivity itself.

\section{Staged experimental program}

\subsection{Design logic}

The program separates six questions:
\begin{enumerate}
  \item Can an explicitly constructed third return state satisfy operational O3 criteria?
  \item Does it preserve directed difference, retained history, and pair specificity?
  \item Can an O3-like contribution arise from ordinary reciprocal learning without an O3 state or objective?
  \item Does the emergent contribution alter both participants and re-enter later intersections?
  \item Is functional return distinct from an external analyst's ability to isolate the contribution?
  \item Is distributed C one jointly nonseparable perspective rather than two additive cross-memories?
\end{enumerate}
Confirmatory thresholds, seeds, controls, and decision rules were frozen before execution. Failed seeds were neither removed nor replaced.

\subsection{Experiments 003R and 004: sufficiency and retained difference}

Experiment 003R corrected an invalid earlier execution and tested a completed explicit C under the full registered control inventory. Among 128 confirmatory pairs, 126 met the registered Phase A Class~2, Phase B Class~2, and Phase B O3 criteria. Across 1,024 Class~0/Class~1 control evaluations, no false Class~2 declaration occurred. A matched generic recurrent control also passed in 126 of 128 pairs; the result supports realizability and transfer, not a unique implementation.

Experiment 004 held $A+B$ constant while reversing $A-B$, matched current inputs while varying earlier interaction history, erased return paths, and exchanged completed carriers between pairs. All four registered hypotheses passed in 128 of 128 pairs. A role-aware recurrent control also passed, identifying a broader architecture class.

\subsection{Experiment 005: spontaneous intersection under ordinary learning}

Two receivers with independent persistent states learned delayed reciprocal recall. After external symbols disappeared, A reported B's earlier symbol and B reported A's. The only training loss was the sum of the two ordinary cross-entropies. There was no C or O3 state, third agent, pair identifier, carrier label, carrier target, or specialized regularizer.

Four equal-capacity interacting architectures and an independent control were evaluated. Post-learning directed intersection contributions were deleted and compared with 64 receiver-norm-matched random directions. A run passed the registered top-0.95 condition when deletion of the extracted contribution exceeded at least the 95th percentile of those controls.

\subsection{Experiments 006 and 006R: self-re-entry}

Experiment 006 froze three consecutive signal-free transitions. At each transition it removed the current contribution, applied one unchanged recurrent update, measured the missing next-state contribution, and reconstructed it. All 96 interaction units transported the contribution above matched random directions at all three transitions. The complete decision nevertheless remained \statusnotsupported\ because the distributed architecture passed a matched individual-history gate in 17 of 24 seeds against the preregistered floor of 18.

Experiment 006R retained that failure and prospectively separated two claims. The primary conjunction tested functional return, bilateral action, cross-pair loss, independent controls, and reconstruction. A separate prediction tested architecture-dependent external identifiability using 48 new seeds per architecture.

\subsection{Experiment 006A: one O3 or two memories}

Experiment 006A tested the decisive ambiguity. Its disconnected dual relay contained two 12-dimensional blocks and exactly the same 486 active parameters, data, optimization budget, and reciprocal-recall objective as each interacting architecture. It could learn both directed recalls but contained no cross-coupling capable of producing a nonadditive joint term.

Four interacting architectures and the disconnected control were evaluated on 48 preregistered seeds, 4,096 held-out episodes, and three signal-free transitions. The ten primary conditions covered capacity, competence, joint norm, bilateral support, learning amplification, transport, bilateral output response, cross-pair exchange, and exact reconstruction. Failure of any condition prevented the primary supported decision.

\section{Results}

\begin{table}[H]
\centering
\footnotesize
\caption{Preregistered evidence from constructed return to emergent operational O3.}
\begin{tabular}{@{}>{\raggedright\arraybackslash}p{1.15cm}>{\raggedright\arraybackslash}p{5.2cm}>{\raggedright\arraybackslash}p{5.0cm}>{\raggedright\arraybackslash}p{2.35cm}@{}}
\toprule
Study & Question & Main observation & Decision \\
\midrule
003R & Can an explicit return structure satisfy operational O3 criteria? & 126/128 pairs; 0/1,024 false Class~2 null calls & Sufficiency supported \\
004 & Does it preserve directed difference, history, transport, and pair specificity? & Four hypotheses passed in 128/128 pairs & Supported \\
005 & Can an intersection-specific contribution emerge without an O3 objective? & 84/96 top-0.95; independent-control norm 0 & Supported \\
006 & Does present C re-enter later states under the original conjunction? & Transport 96/96; distributed separation 17/24 versus floor 18 & \statusnotsupported \\
006R & Can self-re-entry and external identifiability be separated? & 13/13 primary and 3/3 boundary checks passed & Supported \\
006A & Is distributed C one nonseparable O3 rather than two directed memories? & 10/10 checks; interacting C nonzero; competent disconnected C at numerical zero & \statussupported \\
\bottomrule
\end{tabular}
\end{table}

\subsection{Spontaneous formation}

All four Experiment 005 interaction architectures met competence criteria. The extracted directed component exceeded the top-0.95 deletion threshold in 84 of 96 units: 23/24 distributed, 22/24 central shared, 18/24 directional relay, and 21/24 crossbar. Deletion affected both receivers; the independent control produced no corresponding component.

\subsection{Self-re-entry and identifiability}

In Experiment 006, transport passed in all 96 units and exact reconstruction error was below $2.71\times10^{-16}$, but the original conjunction failed as registered. In Experiment 006R, all 13 revised primary conditions passed on unused seeds. All-three-transition transport occurred in 47/48 distributed, 48/48 central-shared, 48/48 directional-relay, and 47/48 crossbar units: 190/192 pooled units against a floor of 168.

Median transported fractions were 0.7941, 0.9336, 0.9207, and 0.9419; median alignments were 0.6487, 0.9039, 0.9209, and 0.9253. Cross-pair exchange increased loss and deletion changed both receivers in every architecture. External-identifiability counts were 29/48, 46/48, 47/48, and 41/48 respectively, supporting all three preregistered architecture-boundary conditions. O3 can therefore be functionally active even when its distributed realization is difficult for an external observer to separate from individual history.

\subsection{006A confirmation of one distributed O3}

Experiment 006A passed all ten primary conditions. All five architectures were competent in 48 of 48 seeds; minimum held-out both-correct accuracy was 0.9751. Median joint-term norms were 0.1555, 0.1303, 0.0936, and 0.1019 for the interacting architectures. The equally competent disconnected dual relay had median norm $2.24\times10^{-17}$ and maximum norm $2.93\times10^{-17}$.

Median bilateral support was 1.0 in every interacting architecture and 0 in the disconnected control. Median transported fractions across three signal-free transitions were 0.5243, 0.5487, 0.4916, and 0.4863, compared with 0.000024 in the disconnected control. Removal changed both receivers' output probabilities in architecture-level medians of 1.0000, 0.9984, 0.9960, and 0.9979 of held-out episodes. Cross-pair substitution imposed positive median cross-entropy cost in all four interacting architectures. Maximum one-step reconstruction error was $5.69\times10^{-16}$.

The 720 transition rows comprise 48 seeds, five architectures, and three transitions. No seed was replaced, no threshold changed, and no NaN or Infinity occurred. Secondary random-direction rankings were not uniformly high; the 006A claim is joint nonseparability against the competent additive control, while Experiments 005 and 006R retain the separately registered random-direction evidence for the full directed carrier.

\section{Discussion}

\subsection{Why C is a third perspective}

C is not called O3 merely because it is a third mathematical term. Under the operational definition, a perspective is a history-bearing organization that transforms received information and constrains subsequent formation. C satisfies that role jointly: it is generated only by reciprocal intersection, has support across both receivers, loses function when either distributed realization is treated separately, changes both participants when removed, is pair-specific, and contributes to the next participant states and next intersection.

The decisive contrast is between competence and intersection. The disconnected relay remembers both directions and solves the same task, but its two memories remain additive and produce numerical-zero joint C. The interacting systems reach comparable competence while producing a nonzero jointly active term. O3 is therefore not equivalent to A's memory of B plus B's memory of A. It is the nonseparable organization of their reciprocal taking-up.

\subsection{Nonlocalization}

No central third object appears in the distributed architecture. One realization of O3 is present in how A has been changed through B; the other is present in how B has been changed through A. Neither contains the whole. Operational unity is established by their inseparable causal function, not by spatial co-location. Experiment 006A supplies the missing test that allows the two distributed realizations to be treated as one O3 rather than named together by convention.

\subsection{Self-re-entry as constraint}

Transport is stronger than persistence of a stored item. Removing $C_t$ changes $h_A(t+1)$ and $h_B(t+1)$ under the same update rule, and changes the next extracted $C_{t+1}$. The present third perspective therefore constrains both the next differentiated perspectives and the next event of intersection. This realizes the computational form of $AB\rightarrow C\rightarrow AB$ as
\[
 (\OA,\OB)\rightarrow\Othree
 \rightarrow(\OA',\OB',\Othree').
\]

\subsection{What changed from relational generation}

Relational generation remains an important consequence, but it is no longer the primary explanatory term. A relation can be externally described without either participant taking up the other. The experiments are more specifically about reciprocal perspective formation. The relation becomes causally generative because intersection produces O3 and O3 re-enters subsequent formation. The conceptual order is therefore perspective difference, reciprocal intersection, emergent O3, and relational self-re-entry.

\subsection{The role of the preserved failure}

Experiment 006 remains negative under its original conjunction. This matters because it prevented functional re-entry from being equated with clean external separability. Experiment 006R then confirmed the separated hypotheses prospectively. The sequence supports a central feature of distributed O3: an operation may be causally coherent across A and B even when an external analyst cannot isolate it uniformly as a dominant direction in each participant.

\section{Subjectivity-Intersection Mathematics as a research program}

The program now has three connected layers:
\begin{enumerate}
  \item \textbf{Perspective mathematics:} formalize differentiated, history-dependent state organizations.
  \item \textbf{Intersection operations:} define reciprocal crossing, nonseparability, identity, exchange, and equivalence.
  \item \textbf{Self-re-entry dynamics:} characterize how generated O3 constrains subsequent perspectives and intersections.
\end{enumerate}

Candidate invariants across the tested implementations are preserved difference, non-reduction to isolated histories, joint nonseparability, pair specificity, bilateral action, signal-free transport, reconstruction through unchanged dynamics, and the distinction between functional activity and external identifiability.

Priority mathematical problems include necessary and sufficient conditions for O3 emergence, equivalence classes of distributed O3 realizations, conservation and transformation laws under repeated intersection, extension beyond two participants, continuous-time intersection, and representation theorems relating localized, shared, and distributed implementations.

\section{Scope and limitations}

The direct result concerns finite, discrete-time recurrent systems trained on synthetic reciprocal-recall tasks. The symbols probe perspective-dependent state formation but do not themselves constitute rich human perspectives. Generalization to biological, social, physical, or phenomenal domains requires domain-specific bridge definitions and interventions.

The operational identification of C with O3 is theoretically stated in Version 3 after the preregistered 006A result. Experiment 006A preregistered joint nonseparability, not the wording of this later theoretical identification. Future tests can preregister the complete O3 definition given in this paper.

The extraction operator in Equation~\eqref{eq:joint} is one preregisterable decomposition. Other operators may reveal equivalent or different intersection structures. The present results establish a reproducible operational object, not uniqueness of decomposition or architecture.

An operational third perspective is not by itself evidence of phenomenal experience. The stronger scientific point is that a third perspective need not be installed as a third agent or localized state: it can be generated, distributed, causally identified, and shown to re-enter the formation of its two participants.

\section{Conclusion}

The staged experiments support a more specific account than relational generation alone. Two independent recurrent perspectives were trained only to answer an ordinary reciprocal task. When each perspective took up information originating from the other, the interacting systems developed a distributed joint contribution absent from an equally competent disconnected pair. The contribution was reducible to neither one-sided history, acted on both receivers, lost function under cross-pair exchange, continued through signal-free transitions, and contributed to the next extracted intersection.

Under the operational definition advanced here, that contribution is a third perspective, O3. It is not a third box between A and B. It is distributed as two internally realized sides of one nonseparable operation: A changed through B and B changed through A. Experiment 006A provides the reason to treat those sides as one O3 rather than two memories, while Experiments 006 and 006R provide the reason to treat O3 as a constraint on subsequent formation rather than a static residue.

The resulting computational statement is concise:
\[
 \boxed{(\OA,\OB)\xrightarrow{\text{intersection}}\Othree
 \xrightarrow{\text{self-re-entry}}(\OA',\OB',\Othree')}
\]
Subjectivity-Intersection Mathematics is therefore proposed as the mathematics of how differentiated perspectives intersect, how a nonseparable third perspective emerges, and how that perspective participates in the next formation of difference and intersection.

\appendix
\section{Reproducibility and public records}

\noindent\textbf{Repository}
\begin{center}
\small\url{https://github.com/SIEL-Research/Subjectivity-Intersection-Mathematics}
\end{center}

\noindent\textbf{Zenodo community}
\begin{center}
\small\url{https://zenodo.org/communities/subjectivity-intersection-mathematics}
\end{center}

\begin{longtable}{@{}>{\raggedright\arraybackslash}p{1.2cm}>{\raggedright\arraybackslash}p{5.0cm}>{\raggedright\arraybackslash}p{7.8cm}@{}}
\toprule
Study & GitHub result implementation & Zenodo result and preregistration \\
\midrule
\endhead
003R & \href{https://github.com/SIEL-Research/Subjectivity-Intersection-Mathematics/tree/main/experiments/003r_subjectivity_agent_class2_o3_corrected}{Corrected operational sufficiency experiment} & Result: \href{https://doi.org/10.5281/zenodo.21761140}{doi:10.5281/zenodo.21761140}; preregistration: \href{https://doi.org/10.5281/zenodo.21761030}{doi:10.5281/zenodo.21761030} \\
004 & \href{https://github.com/SIEL-Research/Subjectivity-Intersection-Mathematics/tree/main/experiments/004_difference_preserving_o3_discrimination}{Difference-preserving and retained-history audit} & Result: \href{https://doi.org/10.5281/zenodo.21762041}{doi:10.5281/zenodo.21762041}; preregistration: \href{https://doi.org/10.5281/zenodo.21761975}{doi:10.5281/zenodo.21761975} \\
005 & \href{https://github.com/SIEL-Research/Subjectivity-Intersection-Mathematics/tree/main/experiments/005_emergent_relational_carrier_solution_class}{Emergent carrier solution class} & Result: \href{https://doi.org/10.5281/zenodo.21763963}{doi:10.5281/zenodo.21763963}; preregistration: \href{https://doi.org/10.5281/zenodo.21763898}{doi:10.5281/zenodo.21763898} \\
006 & \href{https://github.com/SIEL-Research/Subjectivity-Intersection-Mathematics/tree/main/experiments/006_spontaneous_o3_reentry}{Spontaneous operational O3 re-entry} & Result: \href{https://doi.org/10.5281/zenodo.21764472}{doi:10.5281/zenodo.21764472}; preregistration: \href{https://doi.org/10.5281/zenodo.21764327}{doi:10.5281/zenodo.21764327} \\
006R & \href{https://github.com/SIEL-Research/Subjectivity-Intersection-Mathematics/tree/main/experiments/006r_spontaneous_o3_reentry_revised}{Revised confirmation and architecture boundary} & Result: \href{https://doi.org/10.5281/zenodo.21764580}{doi:10.5281/zenodo.21764580}; preregistration: \href{https://doi.org/10.5281/zenodo.21764531}{doi:10.5281/zenodo.21764531} \\
006A & \href{https://github.com/SIEL-Research/Subjectivity-Intersection-Mathematics/tree/main/experiments/006a_emergent_nonseparable_relational_c}{Emergent nonseparable relational C} & Result: \href{https://doi.org/10.5281/zenodo.21785748}{doi:10.5281/zenodo.21785748}; preregistration: \href{https://doi.org/10.5281/zenodo.21785602}{doi:10.5281/zenodo.21785602} \\
\bottomrule
\end{longtable}

The invalid Experiment 003 record remains preserved at \href{https://doi.org/10.5281/zenodo.21760837}{doi:10.5281/zenodo.21760837}; the Experiment 003R technical addendum is \href{https://doi.org/10.5281/zenodo.21761925}{doi:10.5281/zenodo.21761925}. Experiment 006 remains \statusnotsupported. Experiment 006R and 006A are separately preregistered tests on new confirmatory allocations.

\section{Terminology}

\begin{description}[style=nextline,leftmargin=1em,labelindent=0pt]
\item[Operational perspective] A history-dependent state organization that transforms received information and constrains subsequent state and output formation.
\item[Reciprocal perspective intersection] Mutual state formation in which A takes up B through A's organization and B takes up A through B's, while their distinction is preserved.
\item[Directed intersection trace] One receiver-side change attributable to the other under a matched counterfactual.
\item[Operational third perspective O3] A post-learning, jointly nonseparable, pair-specific and self-re-entering contribution generated by reciprocal perspective intersection and distributed across both participants.
\item[Nonseparable C] The inclusion-exclusion term $H(ab)-H(a0)-H(0b)+H(00)$, which vanishes for additive disconnected memories.
\item[Self-re-entry] Contribution of present O3 to the next participant states and next intersection through unchanged dynamics.
\item[External identifiability] An analyst's ability to separate O3 from a matched individual-history contribution; it is distinct from O3's causal function.
\end{description}

\addcontentsline{toc}{section}{Selected references}
\renewcommand{\refname}{Selected references}
\begin{thebibliography}{99}

\bibitem{battaglia2018}
Battaglia, P. W., Hamrick, J. B., Bapst, V., et al. (2018). Relational inductive biases, deep learning, and graph networks. \emph{arXiv:1806.01261}.

\bibitem{kipf2018}
Kipf, T. N., Fetaya, E., Wang, K.-C., Welling, M., and Zemel, R. (2018). Neural relational inference for interacting systems. \emph{Proceedings of ICML}, 2688--2697.

\bibitem{lazaridou2018}
Lazaridou, A., Hermann, K. M., Tuyls, K., and Clark, S. (2018). Emergence of linguistic communication from referential games. \emph{NeurIPS 31}.

\bibitem{locatello2019}
Locatello, F., Bauer, S., Lucic, M., et al. (2019). Challenging common assumptions in unsupervised disentangled representations. \emph{Proceedings of ICML}, 4114--4124.

\bibitem{nosek2018}
Nosek, B. A., Ebersole, C. R., DeHaven, A. C., and Mellor, D. T. (2018). The preregistration revolution. \emph{PNAS}, 115(11), 2600--2606.

\bibitem{rabinowitz2018}
Rabinowitz, N. C., Perbet, F., Song, H. F., et al. (2018). Machine theory of mind. \emph{Proceedings of ICML}, 4218--4227.

\bibitem{rumelhart1986}
Rumelhart, D. E., Hinton, G. E., and Williams, R. J. (1986). Learning representations by back-propagating errors. \emph{Nature}, 323, 533--536.

\bibitem{santoro2017}
Santoro, A., Raposo, D., Barrett, D. G. T., et al. (2017). A simple neural network module for relational reasoning. \emph{NeurIPS 30}.

\bibitem{sukhbaatar2016}
Sukhbaatar, S., Szlam, A., and Fergus, R. (2016). Learning multiagent communication with backpropagation. \emph{NeurIPS 29}.

\bibitem{watanabe2026}
Watanabe, S. (2026). \emph{Subjectivity-Intersection Mathematics: A Mathematical Research Program for Viewpoints, Relational Generation, and Simultaneous Closure}. Research white paper. \url{https://www.researchgate.net/publication/410728192}.

\bibitem{watanabe2026sio}
Watanabe, S. (2026). \emph{Subjectivity Intersection Ontology: A Comparative Ontological Framework for Third Subjectivity}. Version 1.0.0. Zenodo preprint. \href{https://doi.org/10.5281/zenodo.21641878}{doi:10.5281/zenodo.21641878}.

\end{thebibliography}

\vfill
\begin{center}
\small\sffamily
Public Working Preprint v3.0 Draft --- 5 August 2026\\
This public working version records the transition from relational generation to perspective intersection and operational O3 as the central formulation.
\end{center}

\end{document}
